% !TeX root = ../main.tex

\section{Characterization Methodology}
\label{sec:characterization_methodology}

\subsection{Simulation Environment}

The characterization environment is based on the open-source
IIC-OSIC-TOOLS design stack.

The main tools used are:

\begin{itemize}
	\item \code{Xschem} for schematic capture,
	\item \code{ngspice} for electrical simulation,
	\item the SKY130A PDK for transistor models,
	\item scripting for automated data extraction and post-processing.
\end{itemize}


\subsection{Device Testbench}

A dedicated MOS characterization testbench is used to control the
device terminal voltages independently.

The transistor operating point is evaluated while sweeping the relevant
bias quantities such as $V_{GS}$ and $V_{DS}$.

% Insert characterization testbench schematic here.


\subsection{Bias Sweeps}

The characterization is based primarily on DC sweeps.

Typical sweep variables include:

\begin{itemize}
	\item gate-source voltage $V_{GS}$,
	\item drain-source voltage $V_{DS}$,
	\item transistor channel length $L$,
	\item process corner,
	\item temperature.
\end{itemize}


\subsection{Extracted Quantities}

For every operating point, the simulator provides or allows the
derivation of the transistor quantities required for analog design.

The principal metrics are

\begin{equation}
	\frac{g_m}{I_D}
	\label{eq:gmid}
\end{equation}

and

\begin{equation}
	A_{v,\mathrm{int}}
	=
	\frac{g_m}{g_{ds}},
	\label{eq:intrinsic_gain}
\end{equation}

where $\frac{g_m}{I_D}$ represents transconductance efficiency and
$\frac{g_m}{g_{ds}}$ represents the transistor intrinsic voltage gain.


\subsection{Data Processing}

Simulation results are exported and processed to generate reusable
characterization curves and design data.

\begin{designnote}
	
	The characterization scripts and testbenches are kept separate from
	the report itself. The report documents the methodology and the
	resulting design insight, while the repository contains the complete
	implementation.
	
\end{designnote}